\section{Introduction}

On a closed connected spin manifold $M$, index theory for the Dirac operator, in particular the family Atiyah--Singer index theorem is a strong and interesting tool to determine non-trivial topology of the space $\Rpsc(M)$ of Riemannian metrics with positive scalar curvature (psc metrics). However, when the index theoretical information vanishes, it remains a notoriously difficult problem to prove connectedness properties for $\Rpsc(M)$. For $\dim M=3$, a recent result by Bamler and Kleiner~\cite{bamler.kleiner:2019p} building on previous work by Marques~\cite{marques:2012} shows that $\Rpsc(M)$ is either empty or contractible.

The above mentioned applications of index theory rely on the fact that the Dirac operator for a psc metric is an invertible operator. Thus, the space $\Rinv(M)$ of metrics whose Dirac operator is invertible contains $\Rpsc(M)$. As a consequence, non-trivial homotopy groups in~$\Rpsc(M)$, that are detected by (untwisted) index theoretical methods remain non-trivial in the space~$\Rinv(M)$. For example, if the index $\alpha(M)\in \KO_m(\point)$, $m=\dim M$, does not vanish, we have~${\Rpsc(M)\subset\Rinv(M)=\varnothing}$. Similarly, spectral flow of the Dirac operator, or equivalently a family index theorem, can be used to show that $\Rpsc(M)$ and $\Rinv(M)$ have infinitely many components if $m\equiv 3 \mod 4$, $m\geq 7$, and at least two components, if $m\equiv 0,1\mod 8$. Such results use the fact that $\KO_{m+1}(\point)$ is non-trivial in these dimensions~$m$.

However, if the $\KO_{m+1}(\point)$-valued spectral flow between Riemannian metrics $g_1,g_2\!\in\! \Rpsc(M)$ is zero, it is difficult to show that $g_1$ and $g_2$ are in the same connected component of $\Rpsc(M)$. The main result of this article is a proof of the connectedness of $\Rinv(M)$ in the cases $m=2$ and $m=4$. Note that this is the first such connectedness result for $\Rinv(M)$. Our result is in contrast to the setting of metrics with positive scalar curvature in the case $m=4$: Seiberg--Witten theory tells us that on many connected closed oriented $4$-manifolds $\Rpsc(M)$ is not connected \cite[Sections~5 and~6]{ruberman:2001}.

If the index $\alpha(M)$ does not vanish, the space $\Rinv(M)$ is empty, and our setting has to be adapted. Let us recall the Atiyah--Singer index theorem for a closed spin manifold $M$ with its extension using Hitchin's $\alpha$-invariant in $\KO_m(\point)$. Let $\maR(M)$ be the space of Riemannian metrics on $M$.
For any $g\in \maR(M)$ we consider the spinor bundle $\Sigma M=M\times_\rho\Sigma_m$ that arises by considering a real irreducible representations $\rho\colon\Cl_m\to \End(\Sigma_m)$ of the $m$-dimensional real Clifford algebra, $m=\dim M$. Let us summarize some classical facts about such representations, as explained and proven, \eg in \cite[Theorem~I.5.8]{lawson.michelsohn:89}. In dimensions $m\not\equiv 3 \mod 4$, the choice of irreducible representation is unique up to isomorphism (in the sense of real representations).
The representation $\Sigma_m$ carries a complex structure, commuting with the $\Cl_m$-action if ${m\equiv 1,2,3,4,5\mod 8}$, and this complex structure is unique (up to isomorphisms in the sense of complex representations) for $m\equiv 2,4\mod 8$. This implies for $m\equiv 1,2,3,4,5\mod 8$ that~$\Sigma M$ can be obtained from the classical definition of the complex spinor bundle by forgetting the complex structure, and we will identify $\Sigma M$ with the classical spinor bundle in this case. In~the cases $m\equiv 2,3,4\mod 8$, there is even a quaternionic structure on $\Sigma_m$, again unique up to isomorphisms in the sense of quaternionic representations. Thus $\Sigma M$ is a quaternionic vector bundle. We write $\KK\ceq \HH$ for $m\equiv 2,3,4\mod 8$, $\KK\ceq \CC$ for $m\equiv 1,5\mod 8$, and $\KK\ceq \RR$ for~${m\equiv 0,6,7\mod 8}$.

The Dirac operator \smash{$\Dirac^g$} is an unbounded $\KK$-linear operator $\Gamma(\Sigma M)\to \Gamma(\Sigma M)$ with discrete spectrum. Thus all eigenspaces are (sub-)vector spaces over $\KK$.
The $\KO_m(\point)$-valued Atiyah--Singer index theorem says that
\[\alpha(M) =
 \begin{cases}
 \dim_\RR \ker\Dirac^g_+ -\dim_\RR \ker\Dirac^g_-\in \KO_m(\point)\cong \ZZ & \text{if}\ m\equiv0\mod 8,\\
 \dim_\CC \ker\Dirac^g\!\mod 2\in \KO_m(\point)\cong \ZZ/2 & \text{if}\ m\equiv 1\mod 8,\\
 \dim_\HH \ker\Dirac^g\!\mod 2\in \KO_m(\point)\cong \ZZ/2 & \text{if}\ m\equiv 2\mod 8,\\
 \dim_\HH \ker\Dirac^g_+ -\dim_\HH \ker\Dirac^g_-\in \KO_m(\point)\cong \ZZ & \text{if}\ m\equiv 4\mod 8,\\
 0\in \KO_m(\point)\cong \{0\} & \text{if}\ m\equiv 3,5,6,7\mod 8.
 \end{cases}
\]
The value of $\alpha(M)$ does not depend on the Riemannian metric and is thus a ``topological invariant''. More precisely, it only depends on the oriented homeomorphism type of $M$ for ${m\equiv 0,4 \mod 8}$, and it also depends on the differential structure and the spin structure if ${m\equiv 1,2\mod 8}$, but it is always independent of the Riemannian metric $g$.
In terms of the classical $\hat\maA$-genus, we have $\alpha(M)=\hat\maA(M)$ for $m\equiv 0\mod 8$ and $\alpha(M)=2\hat\maA(M)$ for $m\equiv 4\mod 8$.

In all cases, the index theorem provides a lower bound on $\dim_\KK\ker\Dirac^g$.
In fact, let us define
\[|\alpha(M)|\ceq
 \begin{cases}
 \bigl|\dim_\RR \ker\Dirac^g_+ -\dim_\RR \ker\Dirac^g_-\bigr| & \text{if}\ m\equiv0 \mod 8,\\
 1 & \text{if}\ m\equiv 1,2\mod 8 \ \text{and}\ \alpha(M)\neq 0,\\
 0 & \text{if}\ m\equiv 1,2\mod 8 \ \text{and}\ \alpha(M)= 0,\\
 \bigl|\dim_\HH \ker\Dirac^g_+ -\dim_\HH \ker\Dirac^g_-\bigr| & \text{if}\ m\equiv 4\mod 8,\\
 0 & \text{if}\ m\equiv 3,5,6,7,
 \end{cases}
\]
then the Atiyah--Singer index theorem implies $\dim_\KK\ker\Dirac^g\geq|\alpha(M)|$.
We define the space
\[
\Rmin(M) \ceq
\bigl\{g\in \maR(M) \mid \dim_\KK \ker \Dirac^{g}=|\alpha(M)| \bigr\}.
\]
of metrics for which the kernel of the Dirac operator attains this lower bound. In particular $\Rmin(M)=\Rinv(M)$ in the case $\alpha(M)=0$.

From now, one let us assume that~$M$ is connected. Then the set $\Rmin(M)$ is open and dense in $\maR(M)$, see \cite{ammann.dahl.humbert:09}. Metrics in $\Rmin(M)$ are called Dirac-minimal or \Dminimal{}, metrics in the complement $\maR(M) \setminus \Rmin(M)$ are called non-\Dminimal.

The main result of this article is the following theorem.
\begin{Theorem}\label{thm.main}
Let $M$ be a closed connected spin manifold of dimension $m=2$ or $m=4$. Then~$\Rmin(M)$ is connected.
\end{Theorem}

